\makeatletter
\@ifundefined{isChecklistMainFile}{
  % We are compiling a standalone document
  \newif\ifreproStandalone
  \reproStandalonetrue
}{
  % We are being \input into the main paper
  \newif\ifreproStandalone
  \reproStandalonefalse
}
\makeatother

\ifreproStandalone
\documentclass[letterpaper]{article}
\usepackage[submission]{aaai2027}
\setlength{\pdfpagewidth}{8.5in}
\setlength{\pdfpageheight}{11in}
\frenchspacing

\begin{document}
\fi
\setlength{\leftmargini}{20pt}
\makeatletter\def\@listi{\leftmargin\leftmargini \topsep .5em \parsep .5em \itemsep .5em}
\def\@listii{\leftmargin\leftmarginii \labelwidth\leftmarginii \advance\labelwidth-\labelsep \topsep .4em \parsep .4em \itemsep .4em}
\def\@listiii{\leftmargin\leftmarginiii \labelwidth\leftmarginiii \advance\labelwidth-\labelsep \topsep .4em \parsep .4em \itemsep .4em}\makeatother

\setcounter{secnumdepth}{0}
\renewcommand\thesubsection{\arabic{subsection}}
\renewcommand\labelenumi{\thesubsection.\arabic{enumi}}

\newcounter{checksubsection}
\newcounter{checkitem}[checksubsection]

\newcommand{\checksubsection}[1]{%
  \refstepcounter{checksubsection}%
  \paragraph{\arabic{checksubsection}. #1}%
  \setcounter{checkitem}{0}%
}

\newcommand{\checkitem}{%
  \refstepcounter{checkitem}%
  \item[\arabic{checksubsection}.\arabic{checkitem}.]%
}
\newcommand{\question}[2]{\normalcolor\checkitem #1 #2 \color{blue}}
\newcommand{\ifyespoints}[1]{\makebox[0pt][l]{\hspace{-15pt}\normalcolor #1}}

\section*{Reproducibility Checklist}

\vspace{1em}
\hrule
\vspace{1em}

% (AAAI "Instructions for Authors" block removed for the final submission.)

% The questions start here

\checksubsection{General Paper Structure}
\begin{itemize}

\question{Includes a conceptual outline and/or pseudocode description of AI methods introduced}{(yes/partial/no/NA)}
yes --- the differentiable execution (ROM as weights, RAM as a soft tape, control flow as gates) is described conceptually and formally

\question{Clearly delineates statements that are opinions, hypothesis, and speculation from objective facts and results}{(yes/no)}
yes --- proven and measured claims are kept distinct from the explicitly qualitative XAI-feasibility claim

\question{Provides well-marked pedagogical references for less-familiar readers to gain background necessary to replicate the paper}{(yes/no)}
yes --- background on the VCS, the 6502, and ALE is cited

\end{itemize}
\checksubsection{Theoretical Contributions}
\begin{itemize}

\question{Does this paper make theoretical contributions?}{(yes/no)}
yes --- we prove the soft (differentiable) forward execution equals the hard execution bit-for-bit at any finite temperature

	\ifyespoints{\vspace{1.2em}If yes, please address the following points:}
        \begin{itemize}
	
	\question{All assumptions and restrictions are stated clearly and formally}{(yes/partial/no)}
	yes --- the finite-temperature and straight-through assumptions are stated formally

	\question{All novel claims are stated formally (e.g., in theorem statements)}{(yes/partial/no)}
	yes --- the soft-equals-hard equality is given as formal theorems

	\question{Proofs of all novel claims are included}{(yes/partial/no)}
	yes --- full proofs are in the supplementary material

	\question{Proof sketches or intuitions are given for complex and/or novel results}{(yes/partial/no)}
	yes --- proof sketches and intuitions accompany the formal results

	\question{Appropriate citations to theoretical tools used are given}{(yes/partial/no)}
	yes --- citations to the underlying tools (e.g., the straight-through estimator) are given

	\question{All theoretical claims are demonstrated empirically to hold}{(yes/partial/no/NA)}
	yes --- the forward-equivalence theorem is verified quantitatively: all 64 games are byte-identical in RAM and pixel-identical on screen (zero differing bytes and zero differing pixels, measured per byte and per pixel), and the temperature-limit / fully-relaxed regime is measured on the full simulator; the surrogate gradients are additionally illustrated in the qualitative XAI study

	\question{All experimental code used to eliminate or disprove claims is included}{(yes/no/NA)}
	yes --- the conformance harness and the soft-vs-hard verifier are included
	
	\end{itemize}
\end{itemize}

\checksubsection{Dataset Usage}
\begin{itemize}

\question{Does this paper rely on one or more datasets?}{(yes/no)}
yes --- the Atari~2600 cartridge ROMs / ALE games

\ifyespoints{If yes, please address the following points:}
\begin{itemize}

	\question{A motivation is given for why the experiments are conducted on the selected datasets}{(yes/partial/no/NA)}
	yes --- Atari/ALE is the platform on which deep RL was established; all 64 supported games are used

	\question{All novel datasets introduced in this paper are included in a data appendix}{(yes/partial/no/NA)}
	NA --- no novel datasets are introduced

	\question{All novel datasets introduced in this paper will be made publicly available upon publication of the paper with a license that allows free usage for research purposes}{(yes/partial/no/NA)}
	NA --- no novel datasets are introduced

	\question{All datasets drawn from the existing literature (potentially including authors' own previously published work) are accompanied by appropriate citations}{(yes/no/NA)}
	yes --- ALE and the Atari platform are cited

	\question{All datasets drawn from the existing literature (potentially including authors' own previously published work) are publicly available}{(yes/partial/no/NA)}
	yes --- the standard ALE ROM set is publicly obtainable (e.g., via ALE / AutoROM)

	\question{All datasets that are not publicly available are described in detail, with explanation why publicly available alternatives are not scientifically satisficing}{(yes/partial/no/NA)}
	NA --- all datasets used are publicly available

\end{itemize}
\end{itemize}

\checksubsection{Computational Experiments}
\begin{itemize}

\question{Does this paper include computational experiments?}{(yes/no)}
yes --- bit-exact conformance sweeps, throughput benchmarks, and a gradient-based XAI study

\ifyespoints{If yes, please address the following points:}
\begin{itemize}

	\question{This paper states the number and range of values tried per (hyper-) parameter during development of the paper, along with the criterion used for selecting the final parameter setting}{(yes/partial/no/NA)}
	NA --- the paper introduces no learning algorithm with tunable hyperparameters; the emulator is a deterministic re-implementation, so there is no hyperparameter search or final-setting selection (the fixed soft and benchmark settings are listed in the supplementary material)

	\question{Any code required for pre-processing data is included in the appendix}{(yes/partial/no)}
	yes --- the code that prepares each run is included: the deterministic action-stream generator and the standard 60-NOOP$+$4-RESET boot; the ROMs are used unmodified

	\question{All source code required for conducting and analyzing the experiments is included in a code appendix}{(yes/partial/no)}
	yes --- all experiment source is included (attached for review; repository released on acceptance)

	\question{All source code required for conducting and analyzing the experiments will be made publicly available upon publication of the paper with a license that allows free usage for research purposes}{(yes/partial/no)}
	yes --- released publicly on acceptance under the MIT License
        
	\question{All source code implementing new methods have comments detailing the implementation, with references to the paper where each step comes from}{(yes/partial/no)}
	yes --- every differentiable (soft) source in both ports (\texttt{jaxtari/jaxtari/diff/}, \texttt{jutari/src/diff/}) carries module- and function-level comments that name the paper construct each step implements---the five differentiable primitives (one-hot read, softmax dispatch, sigmoid branch, straight-through estimator, distance-softmax relaxed read), the forward-equivalence and temperature-limit theorems, the surrogate-gradient corollary, and the three execution modes---together with the \textsc{xitari} routine each one mirrors

	\question{If an algorithm depends on randomness, then the method used for setting seeds is described in a way sufficient to allow replication of results}{(yes/partial/no/NA)}
	yes --- deterministic action streams and the Mnih-style random-NOOP-reset seeding are described and recorded

	\question{This paper specifies the computing infrastructure used for running experiments (hardware and software), including GPU/CPU models; amount of memory; operating system; names and versions of relevant software libraries and frameworks}{(yes/partial/no)}
	yes --- the supplementary \emph{Throughput} section specifies the CPU (Apple~M1~Max) and GPU (NVIDIA GTX~1080~Ti, 11\,GB, and Quadro RTX~5000, 16\,GB, on an institutional Slurm cluster) models and memory, plus the software stack and versions (JAX~0.10.1 with CUDA~12, Python~3.12, Julia~1.12, Slurm~21.08.5).

	\question{This paper formally describes evaluation metrics used and explains the motivation for choosing these metrics}{(yes/partial/no)}
	yes --- byte-identical RAM / pixel-identical screen (conformance) and env-steps per second (throughput) are defined and motivated

	\question{This paper states the number of algorithm runs used to compute each reported result}{(yes/no)}
	yes --- every reported result states its run count: GPU throughput is measured over ten timed runs per configuration (and the CPU baseline over three), and the conformance/exactness results are deterministic and reported from a single exact run (zero variance by construction)

	\question{Analysis of experiments goes beyond single-dimensional summaries of performance (e.g., average; median) to include measures of variation, confidence, or other distributional information}{(yes/no)}
	yes --- GPU throughput is reported as mean $\pm$ standard deviation over ten timed runs per configuration in the supplementary material's GPU throughput table (the run-to-run spread is below $0.2\%$), and the exactness results are reported as exact distributions (zero variance by construction)

	\question{The significance of any improvement or decrease in performance is judged using appropriate statistical tests (e.g., Wilcoxon signed-rank)}{(yes/partial/no)}
	no --- no significance tests are used---exactness is deterministic and throughput is a benchmark, not a compared-method claim

	\question{This paper lists all final (hyper-)parameters used for each model/algorithm in the paper’s experiments}{(yes/partial/no/NA)}
	yes --- the soft sharpness/temperature, batch sizes, boot sequence, and numeric scope are consolidated in a settings table in the supplementary material


\end{itemize}
\end{itemize}
\ifreproStandalone
\end{document}
\fi